% =====================================================================
%  boxes.tex — صناديق الشرح وأمر البيت
% =====================================================================
\tcbset{
  shatibox/.style={enhanced,breakable,arc=2mm,boxrule=0.6pt,
    left=6pt,right=6pt,top=4pt,bottom=4pt,
    fonttitle=\bfseries,coltitle=white,
    before skip=6pt,after skip=6pt},
}

% ---------- البيت ----------
% \bayt{رقم البيت}{الشطر الأول}{الشطر الثاني}
\newcounter{baytnum}
\newcommand{\bayt}[3]{%
  \begin{tcolorbox}[enhanced,colback=parchment,colframe=gold,boxrule=0.8pt,arc=2mm,
      left=4pt,right=4pt,top=5pt,bottom=5pt,before skip=10pt,after skip=6pt,
      borderline west={2.5pt}{0pt}{gold},borderline east={2.5pt}{0pt}{gold}]
    \large
    \begin{tabularx}{\linewidth}{@{}>{\centering\arraybackslash}m{2.2em}
                                  >{\centering\arraybackslash}X
                                  >{\centering\arraybackslash}X@{}}
      {\normalsize\setcounter{baytnum}{#1}\textcolor{gold}{\textbf{(\thebaytnum)}}} & #2 & #3
    \end{tabularx}
  \end{tcolorbox}\label{bayt:#1}}

% بيت طويل يُكتب شطراه سطرين (عند الحاجة)
\newcommand{\baytlong}[3]{%
  \begin{tcolorbox}[enhanced,colback=parchment,colframe=gold,boxrule=0.8pt,arc=2mm,
      left=4pt,right=4pt,top=5pt,bottom=5pt,before skip=10pt,after skip=6pt,
      borderline west={2.5pt}{0pt}{gold},borderline east={2.5pt}{0pt}{gold}]
    \large
    {\normalsize\setcounter{baytnum}{#1}\textcolor{gold}{\textbf{(\thebaytnum)}}}\quad #2\par
    \hfill #3\par
  \end{tcolorbox}\label{bayt:#1}}

% ---------- معاني الكلمات ----------
\newtcolorbox{mufradat}{shatibox,colback=vocabback,colframe=vocabframe,
  title={معاني الكلمات}}
% \kalima{الكلمة}{معناها}
\newcommand{\kalima}[2]{\textbf{\textcolor{vocabframe}{#1}}: #2\par}

% ---------- المعنى ----------
% \maana  أو  \maana[(الأبيات ١١٣–١١٥)]
\newcommand{\maana}[1][]{\par\smallskip\textbf{\textcolor{maroon}{المعنى\ifx\relax#1\relax\else\ #1\fi: }}}

% ---------- القراءات ----------
\newtcolorbox{qiraat}{shatibox,colback=white,colframe=gold!80!black,
  title={\textcolor{white}{القراءات في البيت}},colbacktitle=gold!80!black}
% \qira{القراء}{القراءة}
\newcommand{\qira}[2]{\noindent\textcolor{gold}{$\triangleleft$}\ #1 \;:\; #2\par}

% ---------- القاعدة ----------
\newtcolorbox{qaida}[1][القاعدة]{shatibox,colback=ruleback,colframe=ruleframe,title={#1}}

% ---------- الأمثلة ----------
\newtcolorbox{amthila}[1][أمثلة من القرآن]{shatibox,colback=exback,colframe=exframe,title={#1}}

% ---------- التحريرات ----------
\newtcolorbox{tahrir}[1][تحريرات (من إتحاف البرية بتحريرات الشاطبية)]{shatibox,
  colback=tahrirback,colframe=tahrirframe,title={#1}}
% بيت من الإتحاف داخل صندوق التحريرات
\newcounter{ithafnum}
\newcommand{\ithaf}[3]{\par{\centering\setcounter{ithafnum}{#1}\textcolor{tahrirframe}{\textbf{(\theithafnum)}}\ #2 \quad\ \ #3\par}\smallskip}

% ---------- تنبيه ----------
\newtcolorbox{tanbih}[1][تنبيه]{shatibox,colback=noteback,colframe=noteframe,title={#1}}

% ---------- الخلاصة ----------
\newtcolorbox{khulasa}[1][خلاصة]{shatibox,colback=sumback,colframe=sumframe,title={#1}}

% فاصل بين الأبيات
\newcommand{\fasl}{\par\smallskip{\centering\textcolor{gold}{\large ۞}\par}\smallskip}
